\section{Collision Checking}

The geometric predicates developed so far can now be combined to answer
questions relevant to robot motion.

\subsection{Segment--Polygon Boundary Intersection}

Given a segment $s$ and polygon $\mathcal{P}$, the segment is tested against
every polygon edge. If
\[
E(\mathcal{P})=\{e_0,e_1,\ldots,e_{n-1}\},
\]
then we evaluate
\[
\operatorname{segmentsIntersect}(s,e_i)
\]
for every edge $e_i$.

The segment intersects the polygon boundary if
\[
\exists e_i\in E(\mathcal{P}) \text{ such that } s\cap e_i\neq\varnothing.
\]
For a polygon containing $n$ edges, this requires $O(n)$ segment-intersection
tests.

\subsection{Why Boundary Intersection Is Not Enough}

Consider a segment whose two endpoints lie completely inside a polygon. The
segment may never cross the polygon boundary. In that case,
\[
\text{boundary intersection}=\text{false},
\]
even though the segment lies entirely inside the obstacle.

This demonstrates an important distinction between \emph{boundary
intersection} and \emph{collision}.

\subsection{Collision-Free Segment}

A segment is considered collision free only if, for every obstacle,
\begin{enumerate}
    \item the segment does not intersect the obstacle boundary,
    \item the first endpoint is not inside the obstacle, and
    \item the second endpoint is not inside the obstacle.
\end{enumerate}

Symbolically,
\[
\operatorname{CollisionFree}(s)
=\neg\bigl(
\operatorname{BoundaryIntersection}(s)
\lor \operatorname{Inside}(s.a)
\lor \operatorname{Inside}(s.b)
\bigr).
\]

Touching an obstacle boundary is currently considered a collision.

\subsection{Collision-Free Path}

A path consisting of waypoints
\[
P_0,P_1,\ldots,P_k
\]
contains the segments
\[
(P_i,P_{i+1}), \qquad 0\leq i<k.
\]

The path is collision free if every one of these segments is collision free:
\[
\operatorname{CollisionFreePath}(\mathcal{P})
=\bigwedge_{i=0}^{k-1}
\operatorname{CollisionFree}\!\left(\overline{P_iP_{i+1}}\right).
\]

This creates the first complete geometric motion-validity test in the
project. Later planners will repeatedly call these collision predicates when
constructing visibility graphs, PRMs, RRTs, and humanoid footstep transitions.
